\section{Taoism}

Taoism is the Chinese tradition rooted in the Tao Te Ching and the Yijing and
developed through later cosmology and inner alchemy. It centers on the Tao, the
nameless source and pattern of all things, and traces how the one unfolds into
the many and how the self returns along the same path. Its method is acting in accord with
the Tao without forcing (wu wei).

\subsection{The Creation Model}

In the beginning there is a \textbf{nameless source}. It cannot be spoken or pointed to. It is silent and empty, formed before there was a sky or a ground, and it stands alone without changing. It is the mother of everything that will come. Out of this source comes the \textbf{one}, a single undivided wholeness, all potential held together in one breath.

The one then divides into \textbf{two poles}, a dark receiving side and a bright active side. The two poles alone would fall back together, so a \textbf{third} arises between them, a blending breath that holds them apart and joins them at once. From this stable trio comes the \textbf{ten thousand things}, the full crowd of the world, each carrying the dark on its back and holding the bright in its arms. The wise do not fight this flow. They act without forcing, they move like water that yields and still wears down stone, and by acting this way they return along the path back to the nameless source.

\subsection{The Way and Its Two Aspects}

The ultimate is the Way (the \textbf{Tao}, also written Dao). It is the origin and the pattern of all things, and it cannot be fully named. The opening of the \textbf{Tao Te Ching} states that the Way that can be told is not the eternal Way.

The Tao shows two aspects that arise together and differ only in name.

\begin{itemize}
 \item Non-being (\textbf{Wu}), the nameless, the origin of heaven and earth.
 \item Being (\textbf{You}), the named, the mother of the ten thousand things.
\end{itemize}

Two further terms mark the threshold of creation.

\begin{itemize}
 \item The Limitless (\textbf{Wuji}), the undifferentiated ground without distinction.
 \item The Supreme Ultimate (\textbf{Taiji}), the first division from which the two forces arise.
\end{itemize}

Later Taoism, especially \textbf{Zhou Dunyi}, binds them in the phrase Wuji and yet Taiji.

\subsection{The Cosmogony of Chapter Forty-Two}

Chapter forty-two of the \textbf{Tao Te Ching} gives the central creation formula. The Tao gives birth to the One. The One gives birth to the Two. The Two gives birth to the Three. The Three gives birth to the ten thousand things. The driving stuff is the vital breath (\textbf{qi}), and the binding agent is the harmonizing breath that blends the poles (\textbf{chongqi}).

\begin{longtable}{l L{8cm} }
\caption{The chapter forty-two sequence read through qi.}\\
\toprule
\textbf{Stage} & \textbf{State of qi} \\
\midrule
Tao & beyond qi, the source \\
One & primordial undifferentiated qi \\
Two & qi split into yin qi and yang qi \\
Three & yin qi, yang qi, and chongqi the harmonizing breath \\
Ten thousand things & qi condensed into all forms \\
\bottomrule
\end{longtable}

\subsection{Yin and Yang}

The two poles are \textbf{yin} and \textbf{yang}. \textbf{Yin} is the receptive, dark, cold, yielding, contracting aspect. \textbf{Yang} is the active, bright, warm, firm, expanding aspect. Neither is pure. Each holds a seed of the other, and each turns into the other at its limit.

\subsection{The Five Phases}

The five dynamic agents (\textbf{wuxing}) are not static elements. They are \textbf{Wood}, \textbf{Fire}, \textbf{Earth}, \textbf{Metal}, and \textbf{Water}. They move through two cycles.

\begin{longtable}{L{4cm} L{4.5cm} L{4.5cm}}
\caption{The generating and overcoming cycles of the five phases.}\\
\toprule
\textbf{Phase} & \textbf{Generating cycle (sheng)} & \textbf{Overcoming cycle (ke)} \\
\midrule
Wood & feeds Fire & parts Earth \\
Fire & forms Earth (ash) & melts Metal \\
Earth & bears Metal (ore) & dams Water \\
Metal & carries Water & chops Wood \\
Water & nourishes Wood & quenches Fire \\
\bottomrule
\end{longtable}

\subsection{The Book of Changes}

The Book of Changes (the \textbf{Yijing}, the I Ching) builds the cosmos by repeated splitting. The Great Treatise states that in the Changes there is the Supreme Ultimate, which generates the two modes, the two modes generate the four images, and the four images generate the eight trigrams. Doubling the trigrams gives the sixty-four hexagrams.

The ladder runs \textbf{Taiji} to the \textbf{two lines} (solid yang and broken yin) to the \textbf{four images} to the \textbf{eight trigrams (bagua)} to the \textbf{sixty-four hexagrams}.

\begin{longtable}{L{2.5cm} L{4cm} L{2.5cm} L{4cm}}
\caption{The eight trigrams.}\\
\toprule
\textbf{Trigram} & \textbf{Lines (bottom to top)} & \textbf{Nature} & \textbf{Attribute} \\
\midrule
Qian & solid, solid, solid & Heaven & strong, creative \\
Kun & broken, broken, broken & Earth & yielding, receptive \\
Zhen & solid, broken, broken & Thunder & arousing, movement \\
Kan & broken, solid, broken & Water & abysmal, danger \\
Gen & broken, broken, solid & Mountain & keeping still \\
Xun & broken, solid, solid & Wind & gentle, penetrating \\
Li & solid, broken, solid & Fire & clinging, light \\
Dui & solid, solid, broken & Lake & joyous, open \\
\bottomrule
\end{longtable}

The trigrams are ordered two ways. The \textbf{Earlier Heaven} arrangement, ascribed to \textbf{Fuxi}, is the ideal symmetric order. The \textbf{Later Heaven} arrangement, ascribed to \textbf{King Wen}, is the order of the changing world used for practical work.

\subsection{The Three Treasures and the Energy Fields}

In Taoist inner alchemy the human is built from three vital energies, the \textbf{three treasures (sanbao)}.

\begin{longtable}{L{3cm} L{5cm} L{4.5cm}}
\caption{The three treasures and the three dantian fields.}\\
\toprule
\textbf{Treasure} & \textbf{Meaning} & \textbf{Energy field (dantian)} \\
\midrule
Jing & essence, the generative root & lower, below the navel \\
Qi & breath, vital energy & middle, center of the chest \\
Shen & spirit, mind, consciousness & upper, between the eyebrows \\
\bottomrule
\end{longtable}

The three energy centers (\textbf{dantian}) are where the treasures are stored and worked. The refinement runs upward. Refine jing into qi, refine qi into shen, refine shen and return to emptiness, and return emptiness to the Tao.

\subsection{The Three Pure Ones}

The highest deities of religious Taoism are the \textbf{Three Pure Ones
(Sanqing)}. They are the personified forms of the primordial breath, dwelling in
the three highest heavens.

\begin{longtable}{L{0.34\linewidth} L{0.48\linewidth}}
\caption{The Three Pure Ones.}\\
\toprule
\textbf{Pure One} & \textbf{Aspect} \\
\midrule
Yuanshi Tianzun & the Celestial Worthy of the Primordial Beginning \\
Lingbao Tianzun & the Celestial Worthy of the Numinous Treasure \\
Daode Tianzun & the Celestial Worthy of the Way and its Virtue, linked to Laozi \\
\bottomrule
\end{longtable}

\subsection{The Eight Immortals}

The \textbf{Eight Immortals (Baxian)} are a group of legendary figures who
attained the Tao and now embody its powers. Each carries an emblem.

\begin{longtable}{L{0.34\linewidth} L{0.48\linewidth}}
\caption{The Eight Immortals.}\\
\toprule
\textbf{Immortal} & \textbf{Emblem} \\
\midrule
Li Tieguai & the iron crutch and gourd \\
Zhongli Quan & the fan that revives the dead \\
Lu Dongbin & the sword \\
Zhang Guolao & the fish drum \\
Han Xiangzi & the flute \\
Cao Guojiu & the jade tablet \\
Lan Caihe & the flower basket \\
He Xiangu & the lotus \\
\bottomrule
\end{longtable}

\subsection{Effortless Action and the Return}

\textbf{Wu wei} is effortless action. It does not mean doing nothing. It means acting in accord with the Tao without forcing, like water that yields and still wears down stone. The sage embodies this. The goal is to \textbf{return to the Tao}, to run the flow backward from the many to the one and rest in the nameless source.

\subsection{Comparison to the Base Models}

Taoism matches the two base models of this collection closely. Like the infinite of Kabbalah and the emptiness of Buddhism, the Tao is a nameless ground beyond name and form, and like both it unfolds the one into the many and calls the self home again. The shared theme is clear. \textbf{The world goes out from a single source and returns to it.}
